\section*{Chapitre \ref{ch:iv1:number-riddles} --- Énigmes sur les nombres}
\markboth{\protect\omTocMark{}{Réponses~: Énigmes sur les nombres}}{\protect\omTocMark{}{Réponses~: Énigmes sur les nombres}}

\begin{solution}{iq:iv1:number-riddles:1}
L'angle est de 78 degrés.

Partons de 12 heures. L'aiguille des minutes parcourt 6 degrés par minute, donc à 36 minutes elle se trouve à 216 degrés. L'aiguille des heures parcourt un demi-degré par minute, donc à 4~h~36 elle se trouve à $4 \times 30 + 36 \times 0.5 = 138$ degrés. L'écart vaut 78 degrés.

\iqwrong{Prendre $120$ degrés pour l'aiguille des heures, comme si elle restait sur le 4 pendant toute l'heure. On obtient alors 96.}
\iqfollow{Quand les aiguilles se superposent-elles pour la première fois après 4~h~00~?}{L'aiguille des minutes gagne $5.5$ degrés par minute, et elle part 120 degrés en retard, donc au bout de $120/5.5 = 21.82$ minutes, soit environ 4~h~21~min~49~s.}
\iqfollow{À quelle fréquence se superposent-elles~?}{Onze fois en 12 heures, donc toutes les $12/11$ d'heure, soit toutes les 65 minutes et environ 27 secondes.}
\end{solution}

\begin{solution}{iq:iv1:number-riddles:2}
$250!$ se termine par 62 zéros.

Un zéro final demande un facteur 10, c'est-à-dire un 2 et un 5. Parmi les facteurs, les 2 sont beaucoup plus nombreux que les 5, donc c'est le nombre de 5 qui décide. Par la formule de Legendre\index{formule de Legendre}, les multiples de 5 donnent $250/5 = 50$, les multiples de 25 donnent 10 de plus, et les multiples de 125 donnent 2 de plus. Le total est 62.

\iqwrong{S'arrêter à $250/5 = 50$. Un multiple de 25 porte deux 5 et un multiple de 125 en porte trois.}
\iqfollow{Quelle factorielle se termine par exactement 5 zéros~?}{Aucune. $24!$ en a 4, et $25!$ en a 6, car le facteur 25 apporte deux 5 d'un coup.}
\iqfollow{Et le dernier chiffre non nul~?}{C'est plus difficile, car il faut suivre les 2 non appariés. Pour $10!$, le nombre vaut 3628800, donc le dernier chiffre non nul est 8.}
\end{solution}

\begin{solution}{iq:iv1:number-riddles:3}
Le dernier chiffre est 9.

Seuls les derniers chiffres comptent pour le dernier chiffre d'un produit. Les puissances de 7 se terminent par 7, 9, 3, 1, puis se répètent avec une période de 4. Comme $2026 = 4 \times 506 + 2$, la puissance a le même dernier chiffre que $7^2 = 49$.

\iqwrong{Utiliser 2026 comme si le dernier chiffre cyclait avec une période de 10. Calculer les premières puissances et lire directement la période.}
\iqfollow{Quel est le dernier chiffre de $7^{7^7}$~?}{Il faut l'exposant modulo 4. Comme $7 \equiv -1$ modulo 4 et que $7^7$ est une puissance impaire, l'exposant vaut $\equiv 3$ modulo 4. Le dernier chiffre est donc celui de $7^3 = 343$, c'est-à-dire 3.}
\iqfollow{Quels chiffres ont une période inférieure à 4~?}{Les chiffres 0, 1, 5 et 6 restent fixes, et 4 et 9 ont une période de 2. Les chiffres 2, 3, 7 et 8 ont une période de 4.}
\end{solution}

\begin{solution}{iq:iv1:number-riddles:4}
31 portes sont ouvertes, et la porte 72 est changée 12 fois.

La porte $n$ est changée une fois pour chaque diviseur de $n$. Les diviseurs vont par paires, $d$ et $n/d$, donc le nombre de changements est pair, sauf si les deux membres d'une paire coïncident, ce qui n'arrive que pour un carré parfait. Les portes ouvertes sont donc les carrés parfaits jusqu'à 1000~: $1^2$ jusqu'à $31^2 = 961$, puisque $32^2 = 1024$.

Pour la porte 72, écrivons $72 = 2^3 \times 3^2$. Un diviseur choisit une puissance de 2 parmi 4 possibilités et une puissance de 3 parmi 3, soit $4 \times 3 = 12$. Comme elle est paire, la porte 72 finit fermée.

\iqwrong{Deviner qu'une porte sur deux finit ouverte, comme si chacune était un pile ou face. Le nombre de changements est entièrement déterminé par les diviseurs de la porte.}
\iqfollow{Quelle porte inférieure à 1000 est changée le plus souvent~?}{La porte 840, qui possède 32 diviseurs, plus que tout autre nombre inférieur à 1000. Elle finit fermée.}
\end{solution}

\begin{solution}{iq:iv1:number-riddles:5}
Écrire deux fois un nombre $n$ donne $1001 \times n$, et $1001 = 7 \times 11 \times 13$.

Écrire $n$ trois rangs à gauche puis ajouter $n$ revient à calculer $1000 n + n$. Par exemple, $123123 = 123 \times 1001$. Diviser par 7, puis par 11, puis par 13 redonne $n$.

\iqfollow{Comment cela donne-t-il un critère de divisibilité par 7~?}{Comme $1000 \equiv -1$ modulo 1001, et donc modulo 7, on découpe le nombre en tranches de trois chiffres à partir de la droite, puis on prend la somme alternée. Pour $123456$, la somme vaut $456 - 123 = 333$, qui laisse le reste 4 dans la division par 7, et $123456$ laisse le même reste. Le même critère fonctionne pour 11 et 13.}
\iqfollow{Que donne l'écriture d'un nombre de deux chiffres deux fois de suite~?}{Cela multiplie par 101, qui est premier. Donc $4747 = 47 \times 101$.}
\iqwrong{Essayer des exemples jusqu'à voir un motif. Les exemples peuvent convaincre, mais la question demande pourquoi, et seul le facteur 1001 l'explique.}
\end{solution}

\begin{solution}{iq:iv1:number-riddles:6}
Le nombre est 2519.

Dans chaque cas, le reste est inférieur d'une unité au diviseur. Donc le nombre plus 1 est divisible par chaque $k$ de 2 à 10, et le plus petit de ces nombres plus 1 est le plus petit commun multiple\index{plus petit commun multiple}~: $\mathrm{ppcm}(2, \dots, 10) = 2^3 \times 3^2 \times 5 \times 7 = 2520$. Le nombre est donc 2519.

\iqwrong{Répondre $10! - 1 = 3\,628\,799$. Cela vérifie bien toutes les conditions, mais ce n'est pas le plus petit~: la factorielle compte des facteurs répétés que le plus petit commun multiple n'utilise pas.}
\iqfollow{Et jusqu'à 20~?}{Le plus petit commun multiple ajoute les nombres premiers 11, 13, 17 et 19 et une puissance supplémentaire de 2~: $2^4 \times 3^2 \times 5 \times 7 \times 11 \times 13 \times 17 \times 19 = 232\,792\,560$. Le nombre est donc $232\,792\,559$.}
\end{solution}

\begin{solution}{iq:iv1:number-riddles:7}
$20!$ se termine par 8 zéros en base 12.

Un zéro final en base 12 est un facteur 12, soit $2^2 \times 3$. Comptons les deux nombres premiers par la formule de Legendre\index{formule de Legendre}. Pour 2~: $10 + 5 + 2 + 1 = 18$. Pour 3~: $6 + 2 = 8$. Chaque facteur 12 utilise deux 2 et un 3, donc les 2 permettent $18/2 = 9$ facteurs 12, et les 3 en permettent 8. C'est le plus petit compte qui l'emporte, donc 8 zéros.

\iqwrong{Se fier à l'habitude de la base 10, qui consiste à ne compter que le plus grand nombre premier. Cela donne 8 ici par hasard, mais il faut comparer le compte de chaque nombre premier à son exposant dans 12, car le nombre premier limitant n'est pas toujours le plus grand. Par exemple, $9!$ contient sept 2 et quatre 3, donc il se termine par 3 zéros en base 12, limité par les 2, et non par 4.}
\iqfollow{Et en base 6 et en base 10~?}{En base 6, chaque zéro demande un 2 et un 3, donc $\min(18, 8) = 8$. En base 10, il faut un 2 et un 5, donc $\min(18, 4) = 4$.}
\end{solution}

\begin{solution}{iq:iv1:number-riddles:8}
$a = 3$ et $b = 2$, donc le nombre est 43272.

Décomposons 72 en ses facteurs premiers entre eux, 8 et 9, et utilisons un critère distinct pour chacun. La divisibilité par 8 ne dépend que des trois derniers chiffres, $27b$. Parmi 270 à 279, seul 272 est un multiple de 8, donc $b = 2$. La divisibilité par 9 dépend de la somme des chiffres, $4 + a + 2 + 7 + 2 = 15 + a$, qui doit être un multiple de 9, donc $a = 3$. Vérification~: $43272 = 72 \times 601$.

\iqwrong{Tester les multiples de 72, dont il y a trop. Utiliser seulement des facteurs premiers entre eux~: 72 vaut $8 \times 9$, et non $6 \times 12$, car un nombre peut être à la fois multiple de 6 et de 12 sans être multiple de 72.}
\iqfollow{La réponse est-elle toujours unique~?}{Non. Pour $7a36b$, $b$ peut valoir 0 ou 8, et ces valeurs donnent 72360 et 73368. Vérifier chaque candidat avant de conclure.}
\end{solution}

\enlargethispage{-3\baselineskip}%
\begin{solution}{iq:iv1:number-riddles:9}
Le reste est 8.

Utilisons les petites puissances. Modulo 17, $3^2 = 9$, $3^4 = 81 \equiv 13 \equiv -4$, et $3^8 \equiv 16 \equiv -1$. Donc $3^{16} \equiv 1$, ce que promet le petit théorème de Fermat\index{petit théorème de Fermat} pour tout nombre non divisible par 17. Or $2026 = 16 \times 126 + 10$, donc
\[
3^{2026} \equiv 3^{10} = 3^8 \times 3^2 \equiv -1 \times 9 = -9 \equiv 8.
\]

\iqwrong{Essayer de réduire l'exposant modulo 17. La période des puissances divise $17 - 1 = 16$, donc on réduit l'exposant modulo 16.}
\iqfollow{Toute base a-t-elle la période 16 ici~?}{Non. Elle divise 16. Modulo 17, 3 a la période 16, mais $2^4 = 16 \equiv -1$, donc 2 a la période 8.}
\iqwhy{L'exponentiation rapide et l'arithmétique modulaire de ce type sont à la base des fonctions de hachage et des générateurs de nombres aléatoires, donc un développeur peut très bien être invité à le coder après l'avoir traité à la main.}
\end{solution}

\begin{solution}{iq:iv1:number-riddles:10}
Le plus petit est 301, et le suivant est 721.

Les cinq premières conditions disent que $n - 1$ est un multiple de $\mathrm{ppcm}(2, 3, 4, 5, 6) = 60$, donc $n = 60k + 1$. Appliquons maintenant la dernière condition, modulo 7. Comme $60 \equiv 4$, il faut $4k + 1 \equiv 0$, ce qui est vrai pour $k = 5$, car $4 \times 5 + 1 = 21$. Donc $n = 301$, soit $7 \times 43$.

Les solutions se répètent tous les $\mathrm{ppcm}(60, 7) = 420$. C'est le théorème des restes chinois\index{théorème chinois des restes}~: avec des modules premiers entre eux, il existe exactement une solution dans chaque période complète. La suivante est donc 721.

\iqwrong{Prendre $2 \times 3 \times 4 \times 5 \times 6 = 720$ comme période, ce qui donne 721 comme plus petit. Le multiple commun est 60, non 720, et 721 n'est que la deuxième solution.}
\iqfollow{Résolvez les restes 2, 3, 2 dans la division par 3, 5, 7.}{Les nombres de reste 2 à la fois pour 3 et pour 7 sont $21k + 2$. Modulo 5, cela donne $k + 2 \equiv 3$, donc $k = 1$, et la réponse est 23.}
\end{solution}

\begin{solution}{iq:iv1:number-riddles:11}
Le plus grand montant que l'on ne peut pas payer est 59, et il existe 30 montants que l'on ne peut pas payer.

Pour deux pièces $a$ et $b$ sans facteur commun, le plus grand montant inatteignable vaut $ab - a - b$, soit $77 - 18 = 59$. Le nombre de montants inatteignables vaut $(a - 1)(b - 1)/2 = 6 \times 10 / 2 = 30$.

Utiliser les restes modulo 7\index{argument des restes}. Utiliser les pièces de 7 seules fait avancer à l'intérieur d'une classe de restes. Il faut donc, dans chaque classe, connaître le plus petit montant atteignable avec les pièces de 11~: ce sont 0, 11, 22, 33, 44, 55 et 66, un pour chaque reste. Une fois une classe atteinte, tout montant ultérieur de cette classe est atteignable. Le dernier montant manqué se situe juste en dessous du plus grand de ces nombres~: $66 - 7 = 59$. Au-delà, chaque classe a commencé.

\iqfollow{Et avec des pièces de 6 et de 9~?}{Leur facteur commun est 3, donc aucun montant qui n'est pas un multiple de 3 n'est jamais atteignable, et il n'existe pas de plus grand montant inatteignable.}
\iqwrong{Multiplier et répondre $77 - 1$, ou penser que 77 est le premier point à partir duquel tous les montants sont atteignables. Les pièces couvrent chaque classe bien avant 77.}
\end{solution}

\begin{solution}{iq:iv1:number-riddles:12}
Vous écrivez le chiffre 1 exactement 301 fois, et les chiffres de 1 à 999 ont pour somme 13\,500.

Considérer les nombres de 0 à 999 comme des chaînes de trois chiffres avec des zéros en tête\index{arguments de symétrie}. Par symétrie, chaque chiffre apparaît aussi souvent à chaque rang. Chaque rang contient chaque chiffre 100 fois, donc le chiffre 1 apparaît $3 \times 100 = 300$ fois. Le nombre 1000 en ajoute une, soit 301.

Pour la somme, chaque rang contribue $100 \times (0 + 1 + \dots + 9) = 4500$, et trois rangs donnent $13\,500$. Les zéros en tête et le nombre 0 n'ajoutent rien à la somme, donc les inclure est sans conséquence.

\iqwrong{Compter combien de nombres contiennent un 1. C'est une autre question~: parmi les nombres de 0 à 999, ceux sans 1 sont au nombre de $9^3 = 729$, donc 271 contiennent un 1. Avec 1000, cela fait 272, alors qu'un nombre comme 111 compte encore trois chiffres dans le décompte de 301.}
\iqfollow{Et le chiffre 0, de 1 à 1000~?}{Ici les zéros en tête comptent vraiment. La réponse est 192~: 9 dans les nombres de 10 à 90, 180 dans les nombres à trois chiffres, et 3 dans 1000.}
\end{solution}

\begin{solution}{iq:iv1:number-riddles:13}
Si l'on peut placer des poids sur les deux plateaux, quatre poids suffisent~: 1, 3, 9 et 27. Avec un seul plateau, il en faut six~: 1, 2, 4, 8, 16 et 32.

Avec les deux plateaux, chaque poids a trois états~: sur le plateau de l'objet, sur l'autre plateau, ou hors de la pesée. Donc $n$ poids donnent $3^n$ sommes signées. Elles sont symétriques par rapport à 0, donc au plus $(3^n - 1)/2$ valeurs positives sont atteignables. Pour $n = 3$, cela fait 13, ce qui est insuffisant, et pour $n = 4$, cela fait 40, exactement assez. Les puissances de 3 y parviennent, car chaque nombre admet une unique écriture en base 3 équilibrée\index{ternaire équilibré}, avec les chiffres $-1$, 0 et 1. Par exemple, $23 = 27 - 3 - 1$, et $40 = 27 + 9 + 3 + 1$.

Avec un seul plateau, chaque poids est présent ou absent, donc il s'agit d'écrire le nombre en binaire\index{bases de numération}. Cinq poids n'atteignent que 31, et six atteignent 63.

\iqwrong{Répondre 1, 2, 4, 8, 16, 32 dans les deux versions. Le binaire est optimal avec un seul plateau, mais il gaspille le troisième état.}
\iqfollow{Quels poids utilise-t-on pour peser 5 grammes~?}{$5 = 9 - 3 - 1$~: on place le 9 sur le plateau opposé à l'objet, et le 3 et le 1 du même côté que l'objet.}
\end{solution}

\begin{solution}{iq:iv1:number-riddles:14}
Les deux derniers chiffres sont 01.

Travailler modulo 100 par élévations au carré successives. $3^4 = 81$, donc $3^5 = 243 \equiv 43$. Puis $3^{10} \equiv 43^2 = 1849 \equiv 49$, et $3^{20} \equiv 49^2 = 2401 \equiv 1$. Les puissances de 3 se répètent donc avec la période 20 modulo 100. Comme $1000 = 20 \times 50$, on a $3^{1000} \equiv 1$.

La raison générale est le théorème d'Euler\index{théorème d'Euler}~: un nombre premier avec 100 a une période qui divise $\varphi(100) = 40$. Ici la période réelle est 20, qui est plus courte.

\iqwrong{Penser que le théorème s'applique à toute base. Il faut que la base n'ait aucun facteur commun avec 100.}
\iqfollow{Quels sont alors les deux derniers chiffres de $2^{1000}$~?}{76. La base 2 a un facteur commun avec 100, donc on découpe 100 en 4 et 25. Modulo 4, la puissance vaut 0. Modulo 25, la période divise 20, donc la puissance vaut 1. Le seul nombre inférieur à 100 qui soit multiple de 4 et qui soit égal à un multiple de 25 plus 1 est 76.}
\end{solution}

\begin{solution}{iq:iv1:number-riddles:15}
La plus haute puissance de 2 est $2^3$, la plus haute puissance de 3 est $3^4$, et 7 ne divise pas le nombre.

Écrire $\binom{100}{50} = 100! / (50!)^2$, et compter les facteurs premiers à l'aide de la formule de Legendre\index{formule de Legendre}. Pour 2~: $100!$ contient $50 + 25 + 12 + 6 + 3 + 1 = 97$, et $50!$ en contient $25 + 12 + 6 + 3 + 1 = 47$, donc la réponse est $97 - 94 = 3$. Pour 3~: $48$ dans $100!$ et $22$ dans $50!$, ce qui laisse $48 - 44 = 4$. Pour 7~: $14 + 2 = 16$ dans $100!$ et $7 + 1 = 8$ dans $50!$, ce qui laisse 0.

Le théorème de Kummer\index{théorème de Kummer} donne un autre point de vue~: la puissance d'un premier $p$ dans $\binom{a + b}{a}$ est égale au nombre de retenues obtenues en additionnant $a$ et $b$ en base $p$. En base 7, 50 s'écrit $101$, et l'ajouter à lui-même donne $202$ sans aucune retenue, donc 7 ne divise pas le nombre. En base 2, $50 = 110010$ et il y a trois retenues.

\iqfollow{Pourquoi le nombre doit-il être pair~?}{$\binom{2n}{n}$ est toujours pair, car $\binom{2n}{n} = 2 \binom{2n - 1}{n - 1}$.}
\iqwrong{Supposer que chaque premier inférieur à 100 doit diviser le nombre, puisque $100!$ les contient tous. Le dénominateur peut annuler chaque occurrence, comme c'est le cas ici pour 7.}
\end{solution}

\enlargethispage{-3\baselineskip}%
\begin{solution}{iq:iv1:number-riddles:16}
Factoriser~: $p^2 - 1 = (p - 1)(p + 1)$, et montrer séparément que 8 et 3 divisent ce produit.

Comme $p$ est impair, $p - 1$ et $p + 1$ sont deux pairs consécutifs. L'un d'eux est un multiple de 4, et l'autre un multiple de 2, donc le produit est un multiple de 8. Ensuite, parmi les trois entiers consécutifs $p - 1$, $p$ et $p + 1$, l'un est un multiple de 3. Ce n'est pas $p$, puisque $p$ est un premier strictement supérieur à 3, donc 3 divise le produit. Comme 8 et 3 sont premiers entre eux, 24 le divise. Pour $p = 5$, $p^2 - 1 = 24$, et pour $p = 7$ il vaut 48, donc 24 ne peut pas être remplacé par un nombre plus grand.

\iqfollow{La primalité est-elle vraiment nécessaire~?}{Non. Il suffit que $p$ soit premier avec 6, donc $25^2 - 1 = 624 = 24 \times 26$ convient aussi.}
\iqwrong{Tester 5, 7 et 11 et déclarer le résultat démontré. Des exemples montrent qu'une affirmation peut être vraie, mais une preuve doit couvrir tous les cas.}
\iqfollow{Existe-t-il un énoncé analogue pour les puissances quatrièmes~?}{Oui. Pour un premier supérieur à 5, $p^4 - 1$ est un multiple de 240, par le même raisonnement avec 16, 3 et 5.}
\end{solution}

\begin{solution}{iq:iv1:number-riddles:17}
Le nombre est 2178, car $4 \times 2178 = 8712$.

L'écrire sous la forme $abcd$, avec $4 \times abcd = dcba$. Le produit a quatre chiffres, donc $a$ vaut 1 ou 2. Le dernier chiffre du produit est $a$, et c'est aussi le dernier chiffre de $4d$, qui est pair, donc $a = 2$. Alors $d \geq 8$, et $4d$ doit se terminer par 2, ce qui exclut $d = 9$ et laisse $d = 8$.

Développons maintenant avec $a = 2$ et $d = 8$~:
\[
4(2008 + 100b + 10c) = 8002 + 100c + 10b.
\]
Cela se simplifie en $13b + 1 = 2c$. Comme $c \leq 9$, $b$ vaut au plus 1. Et $b = 0$ rendrait $2c$ impair, donc $b = 1$ et $c = 7$. Le nombre est donc 2178.

\iqfollow{Quel est le nombre correspondant pour 9~?}{1089, car $9 \times 1089 = 9801$. Placer un 9 au milieu prolonge le motif~: $4 \times 21978 = 87912$, et $9 \times 10989 = 98901$.}
\iqwrong{Chercher parmi des milliers de nombres. Le premier et le dernier chiffre fixent $a$ et $d$ en quelques secondes, et le reste tient en une ligne d'algèbre.}
\end{solution}

\begin{solution}{iq:iv1:number-riddles:18}
La somme vaut 360.

Apparier les nombres par symétrie\index{arguments de symétrie}. Pour $0.abcabc\dots$, prendre son partenaire, dont les chiffres sont $9 - a$, $9 - b$ et $9 - c$. Les chiffres du partenaire sont eux aussi deux à deux distincts. Les deux nombres ont pour somme $0.999\dots = 1$. Aucun nombre n'est son propre partenaire, puisque $a = 9 - a$ n'a pas de solution entière.

Le nombre de tels nombres est $10 \times 9 \times 8 = 720$, ce qui fait 360 paires, chacune de somme 1. La somme totale vaut donc 360.

On peut le vérifier autrement. La valeur de $0.abcabc\dots$ est $abc/999$. Chaque chiffre apparaît à chaque rang 72 fois, et en moyenne la chaîne de trois chiffres vaut $499.5$, donc le total vaut $720 \times 499.5 / 999 = 360$.

\iqwrong{Essayer de sommer dans l'ordre des chiffres. Cela est possible, mais l'appariement supprime tous les calculs. Si une somme présente une symétrie, la chercher d'abord.}
\iqfollow{Qu'en est-il de tous les $0.abab\dots$ avec $a \neq b$~?}{Il y en a 90, et le même appariement avec $9 - a$ donne 45 paires. La somme vaut 45.}
\end{solution}

\begin{solution}{iq:iv1:number-riddles:19}
Il existe 132 tels instants, en 66 paires.

Mesurons le temps $t$ en unités de 12 heures, de sorte que $0 \le t < 1$. L'aiguille des heures se trouve à la fraction $t$ d'un tour complet, et l'aiguille des minutes à $12t$, modulo 1. Une lecture est ambiguë si l'échange des aiguilles donne aussi une heure valide $s$. Cela signifie $s = 12t$ et $t = 12s$, modulo 1. En substituant, $144 t \equiv t$, donc $143 t$ est un entier. Il y a 143 valeurs de $t$ de ce type~: $k/143$ pour $k = 0, \dots, 142$.

Parmi elles, certaines correspondent à des aiguilles confondues, de sorte que l'échange donne la même heure. Ce sont les coïncidences, avec $11t$ entier, et il y en a 11. Elles ne sont pas ambiguës. Il reste donc $143 - 11 = 132$. Chacune est associée à une autre heure qui donne la même image, ce qui forme 66 paires.

\iqwrong{Répondre 143. Les 11 instants où les aiguilles se superposent désignent bien une heure unique, puisque échanger deux aiguilles identiques au même endroit ne change rien.}
\iqfollow{L'instant $t = 1/143$ est-il un exemple réel~?}{Oui. Il s'agit de 5 minutes et environ 2 secondes après midi. L'aiguille des heures est à $1/143$ d'un tour et l'aiguille des minutes à $12/143$. Son échange correspond à $t = 12/143$, soit environ 1 h 00 min 25 s, où l'aiguille des heures est à $12/143$ et l'aiguille des minutes à $144/143 \equiv 1/143$.}
\end{solution}

\begin{solution}{iq:iv1:number-riddles:20}
Le plus grand montant est 43, et tout montant supérieur peut être payé.

Travailler modulo 6\index{argument des restes}. Les boîtes de 6 permettent de progresser à l'intérieur d'une classe de restes, il suffit donc de connaître, dans chaque classe, le plus petit montant atteignable avec les boîtes de 9 et de 20~:
\[
\begin{array}{c|cccccc}
\text{reste} & 0 & 1 & 2 & 3 & 4 & 5 \\ \hline
\text{plus petit montant} & 0 & 49 & 20 & 9 & 40 & 29
\end{array}
\]
Ici $49 = 9 + 40$, $40 = 20 + 20$ et $29 = 9 + 20$. Chacun est le plus petit car atteindre un reste 1, 2, 4 ou 5 exige au moins une boîte de 20, et chaque boîte de 9 supplémentaire ne décale que de 3 modulo 6. Le dernier montant manquant est donc dans la classe de 49, juste en dessous~: $49 - 6 = 43$. Tout montant à partir de 44 se situe au niveau ou au-dessus du plus petit montant de sa classe, donc les boîtes de 6 terminent le travail.

\iqwrong{Appliquer la formule à deux pièces à 9 et 20 et répondre 151. C'est juste si aucune boîte de 6 n'existe. Il n'existe pas de formule aussi simple pour trois valeurs de pièces\index{problème de Frobenius des pièces}, il faut donc construire le tableau.}
\iqfollow{À quoi ressemble 44, pour vérifier~?}{$44 = 20 + 6 \times 4$, et $45 = 9 \times 5$, $46 = 20 + 20 + 6$, $47 = 20 + 9 \times 3$, $48 = 6 \times 8$, $49 = 20 + 20 + 9$. Ces six montants consécutifs, auxquels s'ajoutent les boîtes de 6, couvrent tous les suivants.}
\end{solution}

\begin{solution}{iq:iv1:number-riddles:21}
Les puissances de 2 sont exactement les nombres qui ne peuvent pas s'écrire de cette façon. Et 2026 ne s'écrit que d'une seule manière~: $505 + 506 + 507 + 508$.

Supposons que $k \geq 2$ entiers consécutifs à partir de $m \geq 1$ aient pour somme $n$. Leur somme vaut $k(2m + k - 1)/2$, donc
\[
2n = k \, (2m + k - 1).
\]
Les deux facteurs sont de parités opposées, puisque leur différence $2m - 1$ est impaire. Le second est plus grand que $k$. Si $n$ est une puissance de 2, alors $2n$ l'est aussi, et il n'a aucun diviseur impair supérieur à 1. Le facteur impair parmi les deux doit donc valoir 1. Mais $k \geq 2$ et l'autre facteur est encore plus grand, c'est impossible.

Sinon, $n$ possède un diviseur impair $d > 1$. Écrivons $2n$ comme produit de la partie impaire $d$ et de la partie paire $2n/d$. Appelons $k$ le plus petit des deux et $2m + k - 1$ le plus grand~; ils sont de parités opposées, donc $m$ est un entier, et comme il s'agit du plus grand, $m \ge 1$. Chaque diviseur impair supérieur à 1 donne exactement une façon. Comme $2026 = 2 \times 1013$ avec 1013 premier, il n'y a qu'une façon, avec $k = 4$ et $2m + 3 = 1013$.

\iqfollow{Combien de façons pour 15~?}{Ses diviseurs impairs supérieurs à 1 sont 3, 5 et 15, donc trois façons~: $7 + 8$, $4 + 5 + 6$ et $1 + 2 + 3 + 4 + 5$.}
\end{solution}

\enlargethispage{-3\baselineskip}%
\begin{solution}{iq:iv1:number-riddles:22}
La personne 73 reste.

Commençons par le cas où le nombre de personnes est une puissance de 2. Chaque tour complet du cercle élimine la moitié des personnes et se termine juste avant la personne 1, si bien que le comptage repart de la personne 1. La personne 1 reste donc lorsque le nombre de personnes est une puissance de 2.

Posons maintenant $100 = 64 + 36$. Une fois que 36 personnes sont sorties, il en reste 64, et le comptage est arrivé à la personne qui suit la 36\up{e} éliminée. La 36\up{e} éliminée est la personne 72, donc le prochain survivant est la personne 73, et à partir de là on se retrouve dans le cas de la puissance de 2, avec la personne 73 dans le rôle de la personne 1.

Plus généralement, pour $n = 2^a + b$ avec $0 \le b < 2^a$, le survivant est $2b + 1$\index{problème de Josèphe}. En binaire, cela revient à déplacer le chiffre de tête de $n$ en fin de nombre~: $100$ s'écrit $1100100$, ce qui donne $1001001$, soit 73.

\iqwrong{Répondre 1 par l'argument de la puissance de 2, qui ne tient que lorsque le nombre de personnes est une puissance de 2. Ici la personne 1 survit aux deux premiers tours. Avec 25 restantes, le tour se termine sur une personne qui reste, si bien que la personne 1 sort la première au tour suivant.}
\iqfollow{Et si c'est chaque troisième personne qui sort~?}{La forme simple est perdue. Utiliser la récurrence $J(n) = (J(n - 1) + 3) \bmod n$, en numérotant les personnes à partir de 0. Pour 100 personnes, elle donne le numéro 91, dans la numérotation de cette question.}
\end{solution}

\begin{solution}{iq:iv1:number-riddles:23}
\begin{sloppypar}Le nombre est $105\,263\,157\,894\,736\,842$. Il comporte 18 chiffres, et le doubler donne $210\,526\,315\,789\,473\,684$.\end{sloppypar}

Soit le nombre $N = 10a + d$, où $d$ est le dernier chiffre et $a$ comporte $n - 1$ chiffres. Déplacer $d$ en tête donne $d \cdot 10^{n-1} + a$, et cela doit être égal à $2N = 20a + 2d$. D'où
\[
19a = d \,(10^{n-1} - 2).
\]
Comme $d < 19$ et que 19 est premier, 19 doit diviser $10^{n-1} - 2$. Il faut donc le plus petit $m = n - 1$ tel que $10^m \equiv 2$ modulo 19. En parcourant les puissances de 10, on obtient $m = 17$.

Alors $a = d(10^{17} - 2)/19$. Avec $d = 1$, ce nombre n'a que 16 chiffres, ce qui imposerait un zéro en tête de $N$. La plus petite valeur valable est $d = 2$, qui donne $a = 10\,526\,315\,789\,473\,684$ avec ses 17 chiffres.

On peut aussi le construire à la main en partant de la droite~: on commence par le dernier chiffre 2, et on double à répétition, en écrivant chaque chiffre du produit à gauche du précédent. On s'arrête après avoir écrit un 1 sans retenue, car son double vaut 2, le chiffre de départ. C'est la façon la plus rapide d'obtenir les chiffres à la main.

\iqfollow{Pourquoi ces chiffres sont-ils familiers~?}{Ce sont le bloc périodique de $2/19$ écrit en décimal.}
\iqfollow{Et pour le triplement~?}{Les mêmes étapes avec $29a = d(10^{n-1} - 3)$ donnent une réponse de 28 chiffres, $1\,034\,482\,758\,620\,689\,655\,172\,413\,793$.}
\end{solution}

\begin{solution}{iq:iv1:number-riddles:24}
$C = 7$.

L'astuce consiste à ne pas calculer le nombre. Deux faits suffisent~: la somme des chiffres conserve le reste modulo 9, et les sommes des chiffres diminuent très vite.

D'abord, le reste. $4444 \equiv 16 \equiv 7$ modulo 9. Les puissances de 7 modulo 9 valent 7, 4, 1 puis se répètent, avec une période de 3. Comme $4444 \equiv 1$ modulo 3, $4444^{4444} \equiv 7$ modulo 9, et donc $C \equiv 7$ également\index{somme des chiffres modulo 9}.

Ensuite, la taille. Le nombre est inférieur à $10\,000^{4444} = 10^{17\,776}$, il a donc au plus $17\,776$ chiffres, et $A \le 9 \times 17\,776 = 159\,984$. Un nombre inférieur à cette valeur comporte au plus six chiffres, et le premier de ces six chiffres vaut 1, donc $B \le 1 + 5 \times 9 = 46$. Jusqu'à 46, la plus grande somme de chiffres revient à 39, donc $C \le 12$. Le seul nombre inférieur ou égal à 12 dont le reste modulo 9 vaut 7 est 7.

\iqwrong{Supposer que la somme des chiffres d'une grande puissance vaut environ $4.5$ par chiffre et estimer à partir de là. Cela ne donne $A$ et $B$ qu'approximativement, et c'est le reste qui fixe la réponse.}
\iqfollow{Quelles sont les valeurs exactes~?}{Le calcul direct donne $A = 72\,601$, $B = 16$ et $C = 7$. On remarque que $A$ est bien en dessous de sa borne, près de $4.5$ par chiffre.}
\end{solution}

\begin{solution}{iq:iv1:number-riddles:25}
Le plus petit nombre possible de noix de coco est 3121.

Chaque marin prend un tas $x$ et laisse $\tfrac45(x - 1)$. L'étape astucieuse consiste à ajouter quatre noix de coco fantômes~: comme $\tfrac45(x - 1) + 4 = \tfrac45(x + 4)$, la quantité $x + 4$ se fait simplement multiplier par $\tfrac45$ à chaque étape\index{décalage invariant}. Après cinq marins, on a donc
\[
\text{final} + 4 = \left(\tfrac45\right)^5 (N + 4).
\]
Pour que ce soit entier, $5^5$ doit diviser $N + 4$, donc $N + 4 = 3125\,k$ et le tas final vaut $1024\,k - 4$. Le tas final doit aussi être divisible par 5, et $1024\,k - 4 \equiv 4k - 4$ modulo 5, donc $k \equiv 1$. Le plus petit est $k = 1$, ce qui donne $N = 3121$.

Vérification~: les tas valent $3121$, $2496$, $1996$, $1596$, $1276$, $1020$, et $1020 = 5 \times 204$.

La valeur $-4$ est le point fixe de l'étape, puisque $\tfrac45(-4 - 1) = -4$. Si des noix de coco négatives étaient admises, un tas de $-4$ resterait $-4$ indéfiniment, et c'est ce qui suggère les noix de coco fantômes.

\iqwrong{Raisonner à rebours à partir d'une supposition sur le tas final et essayer chaque valeur à tour de rôle. On finit par trouver 3121, mais il faut essayer beaucoup de valeurs, et on ne peut pas savoir si l'on tient le plus petit.}
\iqfollow{Et si une noix de coco de plus revient au singe le matin~?}{Le tas final doit valoir un de plus qu'un multiple de 5~: $1024\,k - 4 \equiv 1$ impose $k \equiv 0$ modulo 5, donc $k = 5$, ce qui donne $N = 5^6 - 4 = 15\,621$.}
\end{solution}
